\documentclass[12pt]{article}
\usepackage[utf8]{inputenc}
\usepackage[T1]{fontenc}
\usepackage{hyperref}
\title{Tomini Language Resources for K-12 Curriculum Development}
\author{Compiled by Claude Sonnet 5 (claude-sonnet-5) for Andr\'as Kornai}
\date{}
\begin{document}
\maketitle

\section{Introduction}
Tomini (ISO 639-3 txm; also called Tialo) is an Austronesian language of the Celebic branch, spoken along the coast of Tomini Bay in Central Sulawesi, Indonesia. It belongs to the Northern Tomini subgroup with Lauje and Dondo. Wikipedia gives about 30,000 speakers (2001), while Endangered Languages notes that the figures are unreliable and classes the language as threatened.

\section{Primary Educational Resources}
I found no primers, textbooks or school programmes.

\section{Linguistic and Reference Materials}
\begin{itemize}
\item Tomini language (Wikipedia): \url{https://en.wikipedia.org/wiki/Tomini_language}
\item Endangered Languages Project entry: \url{https://endangeredlanguages.com/lang/2793}
\item Verbix, short Tomini verb list with glosses (not a dictionary): \url{https://docs.verbix.com/Languages/Tomini}
\end{itemize}
The search results also pointed to Himmelmann, \emph{Sourcebook on Tomini-Tolitoli Languages: General Information and Word Lists} (Pacific Linguistics 511, 2001). That is the most likely source of word lists, but I did not find a link to it in the results. I found no Tomini--Indonesian dictionary.

\section{Religious and Community Resources}
\begin{itemize}
\item Joshua Project profile of the Tomini people: \url{https://joshuaproject.net/people_groups/15550}
\item Praying for Indonesia, ``Who are the Tomini'': \url{https://www.prayingforindonesia.com/ethnic-groups/the-people-of-sulawesi/who-are-the-tomini/}
\end{itemize}
These are ethnographic and missionary-oriented profiles, not language-teaching resources.

\section{Digital and Audio Resources}
I found none.

\section{Curriculum Development Context}
Tomini is still spoken but is declining among young people, and I found no evidence of school teaching. A curriculum would have to be built from scratch. The Himmelmann sourcebook and fieldwork with the community would be the starting points.
\end{document}